\section{Kết quả thực nghiệm}

Bảng~\ref{tab:final-models} tóm tắt các quyết định cuối cùng. Các mô hình được phân vai như sau:

\begin{itemize}
  \item Conservative full-coverage baseline: \model{text_768_original + minibatch_k16}.
  \item Best experimental full-coverage model: \model{text_pca64_lexical + minibatch_k40}.
  \item Event/dense detection model: \model{text_pca64_only + HDBSCAN(min_cluster_size=30, min_samples=20, cluster_selection_method=leaf)}.
  \item High-silhouette risky reference: \model{text_pca64_only + minibatch_k96}, chỉ dùng để tham chiếu về rủi ro phân mảnh.
\end{itemize}

\begin{table}[H]
\centering
\small
\caption{So sánh rút gọn các mô hình cuối cùng.}
\label{tab:final-models}
\begin{tabularx}{\textwidth}{L{0.24\textwidth} L{0.32\textwidth} r r r r Y}
\toprule
Vai trò & Mô hình & Coverage (\%) & Số cụm & Silhouette & DBI & Ghi chú \\
\midrule
Conservative baseline & \tablemodel{text_768_original + minibatch_k16} & 100.00 & 16 & 0.0632 & 3.0854 & Full-coverage, số cụm vừa phải, dùng làm baseline thận trọng. \\
Selected experimental & \tablemodel{text_pca64_lexical + minibatch_k40} & 100.00 & 40 & 0.1065 & 3.3463 & Cân bằng tốt giữa metric, kích thước cụm và khả năng diễn giải. \\
Event detection & \tablemodel{text_pca64_only + HDBSCAN(mcs=30, ms=20, leaf)} & 20.12 & 219 & 0.6086 & 1.0260 & Phát hiện cụm dày; không nhắm coverage 100\%. \\
Risk reference & \tablemodel{text_pca64_only + minibatch_k96} & 100.00 & 96 & 0.1161 & 2.9968 & Metric mạnh nhưng rủi ro phân mảnh và metadata warning cao. \\
\bottomrule
\end{tabularx}
\end{table}

\begin{table}[H]
\centering
\scriptsize
\caption{Kiểm chứng HDBSCAN cho phát hiện cụm dày/sự kiện.}
\label{tab:hdbscan-event-validation}
\begin{tabularx}{\textwidth}{L{0.28\textwidth} r r r r r r r Y}
\toprule
Cấu hình & Fit rows & Coverage & Noise & Số cụm & Sil. cos & DBI & Neg. sil. & Ghi chú \\
\midrule
\tablemodel{text_pca64_only, mcs=30, ms=20, leaf} & 100000 & 20.12 & 79.88 & 219 & 0.6086 & 1.0260 & 0.69 & Selected event/dense detection model. \\
\tablemodel{text_768_original, hdbscan_minsize50} & 100000 & 20.64 & 79.36 & 105 & 0.5652 & 0.9861 & -- & Historical text-only baseline. \\
\tablemodel{text_pca64_only, mcs=50, eom} & 50000 & 23.20 & 76.80 & 28 & 0.2236 & 1.0837 & 55.09 & Same-sample PCA64 baseline, kém hơn rõ về silhouette. \\
\bottomrule
\end{tabularx}
\end{table}


\subsection{So sánh mở rộng giữa các thuật toán}

Bảng so sánh đầy đủ giữa các mô hình được chọn, các baseline và các cấu hình diagnostic/reference được trình bày trong Phụ lục, Bảng~\ref{tab:extended_algorithm_comparison}. Trong phần chính, báo cáo chỉ giữ lại các nhận xét tổng hợp và một số biểu đồ quan trọng để tránh làm nội dung quá tải.

% Extended algorithm comparison snippet for Overleaf.
\paragraph{So sánh mở rộng các thuật toán.}
Bảng~\ref{tab:extended_algorithm_comparison} tổng hợp thêm các mô hình không được chọn bên cạnh các mô hình cuối cùng. Mục tiêu của bảng không phải là xếp hạng mọi thuật toán bằng một metric duy nhất, mà là làm rõ vai trò của từng nhóm mô hình: full-coverage, event detection và diagnostic/reference. Với những metric không tương thích hoặc không có sẵn, bảng dùng ký hiệu ``--'' để tránh diễn giải quá mức.

\begin{itemize}
    \item \textbf{MiniBatchKMeans} là lựa chọn phù hợp nhất cho full-coverage vì gán cụm cho toàn bộ 100k điểm, chạy ổn định và có thể kiểm soát mức độ phân mảnh bằng tham số $k$. Baseline \texttt{text\_768\_original + minibatch\_k16} giữ vai trò thận trọng, còn \texttt{text\_pca64\_lexical + minibatch\_k40} là mô hình thực nghiệm được chọn nhờ cân bằng tốt hơn giữa silhouette, DBI và kích thước cụm.
    \item \textbf{GMM} là baseline xác suất có ích để đối chiếu với KMeans, nhưng không vượt rõ MiniBatchKMeans trên dữ liệu này. Vì vậy GMM được giữ làm mô hình tham chiếu thay vì mô hình chính.
    \item \textbf{HDBSCAN} không nhằm coverage 100\%, mà dùng để phát hiện cụm dày đặc hoặc cụm dạng sự kiện. Cấu hình \texttt{text\_pca64\_only + HDBSCAN(min\_cluster\_size=30, min\_samples=20, leaf)} được chọn vì silhouette cao, negative silhouette thấp và không có dấu hiệu collapse trong kiểm định full 100k.
    \item \textbf{CLARA} có giá trị như diagnostic baseline medoid-based. Tuy nhiên các kết quả CLARA true vẫn yếu hơn MiniBatchKMeans và chỉ chạy ở mức sample, nên không được chọn làm mô hình cuối cùng.
\end{itemize}

Tóm lại, MiniBatchKMeans được chọn cho nhánh full-coverage vì cân bằng tốt giữa độ phủ, chất lượng cụm và khả năng diễn giải. HDBSCAN được chọn riêng cho event detection vì nó phát hiện các vùng dày đặc thay vì cố gắng gán nhãn cho mọi điểm. GMM và CLARA đóng vai trò baseline/diagnostic, còn các cấu hình high-$k$ như \texttt{minibatch\_k96} chỉ nên xem là tham chiếu vì rủi ro phân mảnh và metadata dominance cao hơn.


\begin{figure}[H]
\centering
\includegraphics[width=0.82\linewidth]{full_coverage_silhouette_comparison.png}
\caption{So sánh \metric{silhouette} của các mô hình full-coverage và reference quan trọng.}
\label{fig:full-coverage-silhouette}
\end{figure}

\begin{figure}[H]
\centering
\includegraphics[width=0.82\linewidth]{full_coverage_dbi_comparison.png}
\caption{So sánh \metric{DBI} của các mô hình full-coverage và reference quan trọng.}
\label{fig:full-coverage-dbi}
\end{figure}

\begin{figure}[H]
\centering
\includegraphics[width=0.78\linewidth]{model_scatter_silhouette_vs_dbi.png}
\caption{Quan hệ giữa \metric{silhouette} và \metric{DBI} theo family thuật toán.}
\label{fig:silhouette-dbi-scatter}
\end{figure}

\begin{figure}[H]
\centering
\includegraphics[width=0.82\linewidth]{event_models_coverage_comparison.png}
\caption{Coverage của các mô hình \algo{HDBSCAN} trong vai trò event/dense detection.}
\label{fig:event-coverage}
\end{figure}

\begin{figure}[H]
\centering
\includegraphics[width=0.82\linewidth]{event_models_cluster_count_comparison.png}
\caption{Số cụm của các mô hình \algo{HDBSCAN} trong vai trò event/dense detection.}
\label{fig:event-clusters}
\end{figure}

\begin{figure}[H]
\centering
\includegraphics[width=0.82\linewidth]{k_sweep_cluster_balance_text_pca64_lexical.png}
\caption{Cân bằng kích thước cụm theo \(k\) cho \model{text_pca64_lexical}.}
\label{fig:k-balance}
\end{figure}

\begin{figure}[H]
\centering
\includegraphics[width=0.82\linewidth]{selected_minibatch_umap_2d.png}
\caption{Trực quan \algo{UMAP} 2D cho mô hình MiniBatch được phân tích sâu. Hình này chỉ dùng như artifact chẩn đoán.}
\label{fig:selected-umap}
\end{figure}
